\documentclass[../cours_1re.tex]{subfiles}
\begin{document}
\setcounter{chapter}{18}
\chapter{Conditionnement, indépendance et Bayes}
\label{ch:19}
\begin{objectif}
Distinguer probabilité d'une intersection, conditionnement et indépendance.
Prérequis : probabilités sur un univers fini, tableaux et arbres de seconde.
\end{objectif}
\section{Changer l'univers de référence}
Sur un univers fini $\Omega$, des poids $p_\omega\ge0$ de somme $1$
définissent $P(A)=\sum_{\omega\in A}p_\omega$.
On retrouve $P(\Omega)=1$, $P(\varnothing)=0$ et, pour des événements disjoints,
$P(A\cup B)=P(A)+P(B)$, car on additionne des termes distincts.
Décomposer $A\cup B$ en $A\setminus B$, $B\setminus A$ et $A\cap B$
prouve $P(A\cup B)=P(A)+P(B)-P(A\cap B)$.
\begin{definition}
Si $P(A)>0$, la probabilité de $B$ sachant $A$ est
$P_A(B)=P(B\mid A)=P(A\cap B)/P(A)$.
\end{definition}
On renormalise les poids dans $A$ pour que leur somme vaille $1$.
Ainsi $P(A\cap B)=P(A)P_A(B)$.
Conditionner n'affirme pas que $A$ cause $B$.

\begin{exemple}
Dans une population fictive de $200$ personnes, $80$ prennent le bus,
$50$ viennent de loin et $30$ cumulent ces deux propriétés.
\begin{center}
\begin{tabular}{lrrr}
\toprule & Loin & Pas loin & Total\\\midrule
Bus &30&50&80\\
Pas de bus&20&100&120\\\midrule
Total&50&150&200\\\bottomrule
\end{tabular}
\end{center}
$P(\text{bus}\cap\text{loin})=30/200$,
$P_{\text{loin}}(\text{bus})=30/50$ et
$P_{\text{bus}}(\text{loin})=30/80$ : trois dénominateurs, trois questions.
\end{exemple}

\section{Partitions et probabilités totales}
\begin{theoreme}
Si $A_1,\ldots,A_r$ sont deux à deux disjoints, réunis égaux à $\Omega$,
et de probabilités strictement positives, alors
\[
P(B)=\sum_{i=1}^r P(A_i)P_{A_i}(B).
\]
\end{theoreme}
\begin{demonstration}
Les événements $B\cap A_i$ partitionnent $B$. L'additivité donne
$P(B)=\sum_iP(B\cap A_i)$, puis on utilise la définition conditionnelle.
Si un $A_i$ a une probabilité nulle, son intersection avec $B$ a aussi
une probabilité nulle ; on omet ce terme sans définir $P_{A_i}(B)$.
\end{demonstration}
Dans un arbre, les probabilités des branches issues d'un même nœud
totalisent $1$. Le produit le long d'un chemin donne la probabilité
de l'intersection ; additionner des chemins disjoints donne celle de leur réunion.
L'ordre de l'arbre est celui du conditionnement, pas nécessairement un ordre temporel.

\section{Indépendance : une propriété à vérifier}
\begin{definition}
$A$ et $B$ sont indépendants si $P(A\cap B)=P(A)P(B)$.
Lorsque $P(A)>0$, cela équivaut à $P_A(B)=P(B)$.
\end{definition}
Cette définition reste valable si l'un des événements est de probabilité nulle.
Elle est symétrique ; le conditionnement seul ne l'est pas.
\begin{demonstration}[indépendance des complémentaires]
Si $A$ et $B$ sont indépendants,
\[
P(A\cap\overline B)=P(A)-P(A\cap B)
=P(A)(1-P(B))=P(A)P(\overline B).
\]
On peut échanger les rôles puis remplacer les deux événements par leurs
complémentaires. Toutes ces paires sont indépendantes.
\end{demonstration}
\begin{attention}
Des événements incompatibles ont une intersection vide. S'ils ont tous deux
une probabilité positive, ils ne sont pas indépendants :
$0=P(A\cap B)\ne P(A)P(B)$.
Des observations apparemment sans lien ne suffisent pas à prouver une indépendance.
\end{attention}

\section{Inverser un conditionnement : formule de Bayes}
\approfondissement
Si $P(A)>0$ et $P(B)>0$, la même intersection calculée de deux façons donne
\[
P_B(A)=\frac{P(A)P_A(B)}{P(B)}.
\]
Avec une partition, remplacer le dénominateur par la formule totale.
\begin{exemple}[modèle fictif de contrôle]
Une pièce est défectueuse avec probabilité $0{,}01$. Un contrôle est positif
avec probabilité $0{,}9$ sur une pièce défectueuse et $0{,}05$ sur une pièce
non défectueuse. Ces nombres sont inventés pour le calcul.
En notant $D$ le défaut et $+$ le résultat positif :
\[
P(+)=0{,}01\times0{,}9+0{,}99\times0{,}05=0{,}0585,
\qquad
P_+(D)=\frac{0{,}009}{0{,}0585}=\frac2{13}\simeq15{,}38\%.
\]
Sur $20\,000$ pièces suivant exactement ces proportions, on attend
$180$ positifs défectueux et $990$ positifs non défectueux.
Le faible taux initial de défaut explique le résultat.
\end{exemple}

\begin{finchapitre}
\section{Utilité et prolongements}
Le conditionnement organise les informations disponibles et les arbres
permettent de décomposer une situation complexe.

Bayes intervient en inférence statistique ; l'espérance conditionnelle du
chapitre 28 remplacera un événement par une quantité à prédire.
Les processus stochastiques de licence et de master décriront l'évolution
de ces informations au cours du temps.
\begin{retenir}
Toujours identifier l'univers de référence. Une indépendance est une
égalité, une partition une décomposition, et un conditionnement nécessite un dénominateur non nul.
\end{retenir}
\end{finchapitre}
\end{document}
